
\subsection{Round six: discussion}
\label{sec:r6-discussion}
The one-sentence answer, built from the literal statuses of H33 to H37: At no non-leaky alignment tested (a in {0, 0.25, 0.5, 0.75}) does the literal H34 rule support a wirelength gain of HDL-aware packing over stock VPR (statuses: a=0 not supported, a=0.25 not supported, a=0.5 not supported, a=0.75 not supported); with the leak (a=1) the wirelength ratio is 0.669 (supported); it hurts (WL at least 2\% worse, Holm) at a in {0, 0.25, 0.5, 0.75}. The manipulation did what it was built to do (H30 supported: Spearman 0.988, median measured purity 0.100 at $a=0$ and 1.000 at $a=1$), so a level is a measured fraction of the word file that equals the block partition. The dose-response slope of the wirelength ratio on measured alignment is -1.317 (status supported); a gain concentrated at $a=1$ would also give a negative slope, so the leak-free variant and the per-level tests H34 carry the non-leaky claim. On the chain control, where round five found no headroom, norns / stock is 4.150 (H37 not supported). The placebo (stock at level $a$ against stock at $a=1$; same logic, different register names) has geometric means 1.000 at $a=0$, 1.000 at $a=0.25$, 1.000 at $a=0.5$, 1.000 at $a=0.75$, the size of the movement that renaming alone produces, against which a Norns effect is called found only if it lies outside the placebo interval. Limits: synthetic designs from a generator; the alignment is a construction (a bit-slice deal of flip-flop names), not real HDL, where a word's members need not be bit-slices and LUT words exist; the word file carries flip-flop words only; the Norns packer was not tuned for partial alignment; three topologies and one logic family; the known optimum is a clustering optimum, not a proven wirelength or delay optimum. A result at a given $a$ says what this packer does when that fraction of the HDL words falls in the block, not what real HDL gives.
